% 概要
\section{この資料について}

\subsection{位置づけ}

本リポジトリには3つのドキュメントがある。

\begin{tabularx}{\textwidth}{l l X}
\toprule
\textbf{ドキュメント} & \textbf{ディレクトリ} & \textbf{内容} \\
\midrule
設計仕様書 & \texttt{doc/} & アーキテクチャのルール・API仕様（「何を守るか」） \\
教科書 & \texttt{doc-textbook/} & パターンの段階的な学習（「なぜそうするか」） \\
\textbf{実装ガイド}（本書） & \texttt{doc-guide/} & 実コードの解説（「実際にどう書いたか」） \\
\bottomrule
\end{tabularx}

\begin{cautionbox}[サンプルコードについて]
本リポジトリに実装されているモジュール群のコードは、あくまで\textbf{アーキテクチャの設計パターンを具体的に示すためのサンプル実装}である。
実機での動作検証は行っていないため、実プロジェクトへの適用時は十分なテストを実施すること。
\end{cautionbox}

\subsection{読み方}

本書では、リポジトリに実装されている各モジュールを取り上げ、以下の構成で解説する。

\begin{enumerate}
    \item \textbf{概要} — モジュールの役割と、そのモジュールで学べるパターン
    \item \textbf{コード解説} — 実際のコードを引用し、設計ルールとの対応を示す
    \item \textbf{ポイント} — 特に注目すべき設計判断やテクニック
\end{enumerate}

モジュールは簡単なものから順に並べている。
前半のモジュールで基本パターンを押さえたうえで、後半の応用パターンに進む構成になっている。

\subsection{モジュール一覧}

\begin{longtable}{l l l}
\toprule
\textbf{モジュール} & \textbf{種別} & \textbf{主な学習パターン} \\
\midrule
\endhead
LedModule & 出力 & 最小構成の基本パターン \\
ButtonModule & 入力 & デバウンス（ModuleTimer活用） \\
BatteryModule & 入力 & ADC読み取り + 移動平均フィルタ \\
ServoModule & 出力 & PWM制御（LEDC） \\
BuzzerModule & 出力 & deinit()によるリソース解放 \\
Mpu6500Module & 入力 & I2Cバスパターン + レジスタ直接アクセス \\
TouchModule & 入力 & バスオブジェクト共有（TFT\_eSPIドライバ） \\
TftModule & 出力 & SPI共有 + 複合デバイスの分離 \\
CameraModule & 入力 & PSRAMバッファ + フレーム管理 \\
DriveMotorModule & 出力 & PWM 2ピン制御 + 統合モジュール部品 \\
ChassisModule & 出力 & 機能統合系（内部にモジュールを持つ） \\
BleModule & 入力 & スレッドセーフティ（volatileフラグ方式） \\
WifiModule & 入力 & ステートマシン内蔵 \\
EncoderModule & 入力 & ハードウェア割り込み連携 \\
\bottomrule
\end{longtable}
